\documentclass[11pt,a4paper]{article}

\usepackage[utf8]{inputenc}
\usepackage[T1]{fontenc}
\usepackage[english]{babel}
\usepackage{kpfonts}
\usepackage[
  a4paper,
  top=2.54cm,
  bottom=2.54cm,
  left=3.17cm,
  right=3.17cm]{geometry}
\usepackage{graphicx}
\usepackage{fancyhdr}
\usepackage{setspace}
\usepackage{xcolor}
\usepackage{lipsum}
\usepackage{hyperref}
\usepackage{array}

% Define main color
\definecolor{maincolor}{RGB}{144,41,42}

% Hyperlink colors
\renewcommand\UrlFont{\color{maincolor}\ttfamily}

% Set page style to empty globally (no header or footer)
\pagestyle{empty}

% Set header and footer
\setlength{\headheight}{30pt}
\setlength{\headsep}{21pt}
\setlength{\footskip}{30pt}

% Title formatting
\makeatletter

% Redefine \bfseries to include main color
\DeclareRobustCommand{\bfseries}{%
	\not@math@alphabet\bfseries\mathbf%
	\fontseries\bfdefault\selectfont%
	\color{maincolor}%
}%

% Define title components
\def\@maintitle{}
\def\@subject{}
\def\@subtitle{}
\def\@authorname{}
\def\@delivery{}

% Commands to set title components
\newcommand{\maintitle}[1]{\def\@maintitle{#1}}
\newcommand{\subject}[1]{\def\@subject{#1}}
\newcommand{\subtitle}[1]{\def\@subtitle{#1}}
\newcommand{\authorname}[1]{\def\@authorname{#1}}
\newcommand{\delivery}[1]{\def\@delivery{#1}}

% Redefine \maketitle
\renewcommand{\maketitle}{
    \thispagestyle{empty}
    \begin{flushright}
        {\fontsize{28pt}{32pt}\selectfont \bfseries \color{maincolor} \@maintitle}\\[0.35cm]
        {\fontsize{14pt}{16pt}\selectfont \bfseries \color{maincolor}
        \textit{\@subtitle}}
    \end{flushright}
    \vspace{1cm}
    {\noindent
        \fontsize{12pt}{14pt}\selectfont
        \setlength{\baselineskip}{14pt} 
        \begin{tabular}{@{}l >{\raggedright\arraybackslash}p{0.8\linewidth}}
            \textbf{\textcolor{maincolor}{Subject:}}  & \@subject    \\[6pt]
            \textbf{\textcolor{maincolor}{Authors:}}  & \@authorname \\[6pt] 
            \textbf{\textcolor{maincolor}{Group:}} & T   \\
            \textbf{\textcolor{maincolor}{Delivery:}} & \@delivery   \\
        \end{tabular}
    }
    \vspace{0.5cm} \\
    \noindent\textcolor{maincolor}{\rule{\textwidth}{2pt}}
}
\makeatother

% Set title information
\maintitle{A3 Project Proposal}
\subtitle{Master in Health Data Science --- MHEDAS}
\subject{Complex Networks}
\authorname{Teferi Mekonnen Yitayew, Sofía González Estrada, Ravneet-Rahul Sandhu Singh}
\delivery{December 14, 2025}

% Paragraph formatting
\setlength{\parindent}{0pt}
\onehalfspacing

\begin{document}

\maketitle

\vspace{10pt}

\section*{Option 1:}

\textbf{Title:} The Architecture of Guardiola’s Tiki-Taka: A Complex Network
Analysis of the 2011 UCL Final

\vspace{1em}

\textbf{Description:} The 2011 Champions League Final between FC Barcelona and
Manchester United is widely regarded as a masterpiece of tactical football,
showcasing the pinnacle of Guardiola's tiki-taka philosophy. This project aims
to study the match as a Complex Network by modeling player interactions as a
directed graph, where nodes represent players and edges represent passes. We
aim to quantify the structural properties that defined Guardiola’s dominance.
We hypothesize that his system exhibited specific topological features: high
clustering coefficients and robust small-world properties that distinguished it
from the inefficient structure of Manchester United. Furthermore, we aim to
analyze the specific role of Lionel Messi as a False Nine hub, whose
positioning and passing patterns were crucial to unlocking defenses.

\vspace{1em}

\textbf{Data:} To conduct this analysis, we will utilize event data from the
2011 Champions League Final, sourced from StatsBomb. This dataset provides
event-level information, including pass locations, timestamps, and player
identities.

\vfill

\hfill
\begin{minipage}[b]{0.3\textwidth}
	\includegraphics[width=\linewidth]{../../images/logos/mhedas_logo.png}
\end{minipage}

\cleardoublepage

\section*{Option 2:}

\textbf{Title:} Sophie's Complex World: A Network Analysis of Gaarder's
Philosophical Novel

\vspace{1em}

\textbf{Description:} \textit{El Mundo de Sofía} (Sophie's World) is a novel
written by Jostein Gaarder, who describes the history of philosophy through the
narration of a teenage girl's life. We aim to describe the story from a complex
network perspective, representing each character as a node, with edges between
them defined if two characters are mentioned within a predefined word window.
By analyzing the network, we would try to describe communities formed by
different philosophers. We also want to identify important individuals, like
Alberto Knox, who acts as the guide throughout the story.

\vspace{1em}

\textbf{Data:} To conduct this analysis, we will utilize the full digital text
of \textit{El Mundo de Sofía} in plain text format. We will employ Natural
Language Processing (NLP) techniques to extract entity interactions and build
the adjacency matrix for the network analysis.

\section*{Option 3:}

\textbf{Title:} Mapping the Multiverse: Community Detection in the Marvel
Social Network

\vspace{1em}

\textbf{Description:} The Marvel Universe consists of a vast interconnected
Multiverse of narratives consisting of over 8,000 characters and decades of
crossovers. The idea would be about analyzing this vast lore as a social
collaboration network. We will construct a graph where nodes represent
characters and edges are defined by co-occurrence in the same comic. Our goal
is to map the topology of this Multiverse using Centrality Measures (Degree and
Betweenness) to determine if icons like Captain America are truly the
structural hubs (like in the movies), or if characters act as the primary
bridges between different storylines (e.g., connecting the Avengers cluster to
the X-Men cluster). We will use Community Detection algorithms to see if the
structure of the graph naturally reproduces the teams without prior knowledge
of the lore.

\vspace{1em}

\textbf{Data:} We will utilize the open-source Marvel Universe Social Network
dataset, which maps character appearances across decades of publications.

\vfill

\hfill
\begin{minipage}[b]{0.3\textwidth}
	\includegraphics[width=\linewidth]{../../images/logos/mhedas_logo.png}
\end{minipage}

\end{document}
